\section*{LỜI MỞ ĐẦU}
\addcontentsline{toc}{section}{LỜI MỞ ĐẦU}

\subsection*{1. Tính cấp thiết của đề tài}
Trong bối cảnh nền kinh tế số và cạnh tranh toàn cầu, quản lý và quản trị doanh nghiệp đóng vai trò quyết định đến sự tồn tại và phát triển bền vững của tổ chức. Tập đoàn FPT là một trong những doanh nghiệp công nghệ thông tin và viễn thông hàng đầu Việt Nam, sở hữu mô hình quản lý tập đoàn đa ngành với quy mô nhân sự lớn. Việc nghiên cứu thực trạng quản lý, quản trị tại FPT cùng tác động của Trí tuệ Nhân tạo (AI) mang tính thời sự và có ý nghĩa thực tiễn sâu sắc.

\subsection*{2. Mục tiêu nghiên cứu}
\begin{itemize}
    \item Hệ thống hóa các khái niệm cơ bản về khoa học quản lý và quản trị doanh nghiệp.
    \item Phân tích thực trạng bộ máy quản lý và cơ cấu quản trị tại Tập đoàn FPT.
    \item Đánh giá các điểm nghẽn, hạn chế trong công tác quản lý và đề xuất giải pháp cải thiện.
    \item Đánh giá tác động đa chiều của AI đối với hoạt động quản lý vận hành và quản trị chiến lược tại FPT.
\end{itemize}

\subsection*{3. Đối tượng và phạm vi nghiên cứu}
\begin{itemize}
    \item \textbf{Đối tượng:} Cơ cấu tổ chức bộ máy quản lý, hệ thống quản trị công ty và ứng dụng AI tại Tập đoàn FPT.
    \item \textbf{Phạm vi:} Dữ liệu hoạt động giai đoạn 2020--2024 tại Công ty Cổ phần FPT và các công ty thành viên.
\end{itemize}

\subsection*{4. Phương pháp nghiên cứu}
Phương pháp thu thập và phân tích dữ liệu thứ cấp từ báo cáo thường niên, báo cáo quản trị, phương pháp so sánh đối chiếu và tổng hợp lý luận.
